% docs/thesis/tab-retrieval-leakage-two-subjects.tex
% Measured 2026-08-09 | apd 187 materials / 1193 chunks, gold evals/apd/gold.jsonl (48 questions)
%                     | passc 15 materials / 341 chunks, gold evals/passc/gold.jsonl (76 questions)
% Source: the `at_k` block, field `leakage`, of
%   apd   results-gptoss-glean0.json (ts 2026-08-09T08:57:48Z)
%   passc evals/passc/results-20260809T111747+0000.json (ts 2026-08-09T11:17:47Z)
%   divided by each run's recorded leakage_base_rate:
%   apd   0.037720033528918694 = 45 of 1193 chunks (Examen.md 24, Quiz 1 - ThreadsSynchro.pdf 10,
%         REZUMAT.md 8, Quiz 2 - OpenMP.pdf 3), cross-checked against store.db
%   passc 0.014662756598240470 =  5 of  341 chunks (subiecte.md), cross-checked against store.db
% Config: both subjects on their own openai/gpt-oss-120b, graph.gleanings=0,
%   graph.context_window=16384 graph; rerank=true; arms vector + hybrid-local; provider-free.
% WHY ENRICHMENT, NOT THE RAW RATE: question-source files are themselves indexed, so a raw leakage
%   rate rewards an arm for returning more of the corpus, and the two subjects' base rates differ
%   by 2.6x so raw rates are not comparable across subjects at all. Reported as
%   leakage / leakage_base_rate.
% Matched cutoffs only (median_retrieved_n == k in every cell); errors = 0 in both runs.
% Raw rank-1 leakage rates behind the k=1 column, for the "roughly half" claim in the caption:
%   apd vector 0.5625, apd hybrid-local 0.2708; passc vector 0.4605, passc hybrid-local 0.2895.
% NOT a fourth cell of the apd three-build experiment: Table~\ref{tab:retrieval-leakage} varies the
%   graph on one corpus; this table varies the corpus. The apd columns here are that table's
%   vector and gptoss@0 columns.
\begin{table}[htbp]
  \centering
  \caption{Question-source leakage as enrichment over the corpus base rate, both pilot subjects
    (apd $n=48$, passc $n=76$, 0 errors). A value of 1.0 would mean the returned set contains
    question-source chunks at exactly the rate they occur in the corpus. Contamination is
    concentrated at rank 1 on both subjects and is far more extreme on passc, whose base rate is
    2.6 times lower: \texttt{vector}'s first hit is question-source material 56.3\% of the time on
    apd and 46.1\% on passc, so on both corpora roughly half of all rank-1 chunks are the question
    rather than the material answering it. \texttt{hybrid-local} is consistently the less
    contaminated arm at $k=1$ on both subjects.}
  \label{tab:retrieval-leakage-two-subjects}
  \begin{tabular}{rrrrr}
    \toprule
    & \multicolumn{2}{c}{apd (base 3.77\%)} & \multicolumn{2}{c}{passc (base 1.47\%)} \\
    \cmidrule(lr){2-3} \cmidrule(lr){4-5}
    $k$ & \texttt{vector} & \texttt{hybrid-local} & \texttt{vector} & \texttt{hybrid-local} \\
    \midrule
     1 & 14.91 &  7.18 & 31.41 & 19.74 \\
     8 &  3.59 &  3.45 &  7.29 &  7.29 \\
    20 &  2.54 &  2.26 &  5.88 &  4.98 \\
    \bottomrule
  \end{tabular}

  \vspace{0.5em}
  \footnotesize
  Enrichment factor $=$ measured leakage rate $/$ corpus base rate. The base rates recorded in the
  result files are $0.0377$ on apd (45 of 1193 chunks, across four indexed exam and quiz files) and
  $0.0147$ on passc (5 of 341 chunks, all from \texttt{subiecte.md}). Raw rates are not quoted:
  because question-source files are part of the indexed corpus, a raw rate flatters whichever arm
  returns the larger share of the corpus, and the two subjects' base rates differ enough that raw
  rates would not be comparable across them either. Gold \texttt{expected} labels may never name a
  question-source file --- \texttt{eval/gold.py} rejects that at load --- so this contamination
  never inflates a hit directly; it does mean the absolute hit rates in
  Table~\ref{tab:retrieval-quality-two-subjects} are optimistic on both subjects, and it is the
  argument for a contamination-controlled re-index. It does not affect the arm comparison: both
  arms face the same corpus.
\end{table}
